\documentclass[10pt,a4paper]{article}
\usepackage[T1]{fontenc}
\usepackage[utf8]{inputenc}
\usepackage{lmodern}
\usepackage{amsmath,amssymb}
\usepackage[margin=21mm]{geometry}
\usepackage{microtype}
\usepackage[hidelinks]{hyperref}
\usepackage{enumitem}
\setlist[itemize]{leftmargin=*,nosep,topsep=3pt}
\setlength{\parindent}{0pt}
\setlength{\parskip}{5pt}
\setlength{\emergencystretch}{3em}
\newcommand{\question}[1]{\par\medskip{\normalsize\bfseries #1}\par\smallskip}
\pagestyle{plain}

\begin{document}
{\Large\bfseries Pathological pallidal beta activity in Parkinson's disease is sustained during sleep and associated with sleep disturbance\par}

\textbf{Paper:} Yin et al. (2023).\quad Human overnight recordings; offline controller simulation.\par
\href{https://doi.org/10.1038/s41467-023-41128-6}{DOI: \nolinkurl{10.1038/s41467-023-41128-6}}\par
{\footnotesize Page references below use PDF page numbers, including front matter.}

\question{To which general question does this paper relate?}
Measurement and sleep-related control objectives. Tests whether a waking motor biomarker can be interpreted the same way during sleep.

\question{How does it answer the general questions?}
\textbf{Measured.} Pallidal LFPs were recorded alongside polysomnography in patients off dopaminergic medication, with a dystonia comparison group. Parkinsonian beta remained elevated during wakefulness, NREM and REM sleep and was associated with disturbed sleep (PDF pp.~1--7).

\textbf{Usable insight.} Offline adaptive-DBS simulations using waking beta thresholds produced sleep-dependent stimulation patterns. A controller calibrated during wakefulness can therefore misrepresent therapeutic demand during sleep. These simulations do not demonstrate improved or worsened sleep under an actual adaptive intervention (PDF pp.~7--9).

\textbf{Model relevance.} Supports including operating context and evaluating sleep outcomes separately from motor scores. The results are correlations; suppressing nocturnal beta is not established as a sufficient treatment. No delayed neural model is identified, and the target is pallidal: its biomarker behaviour should not automatically be assigned to STN recordings.

\question{What papers does it cite that could provide more information?}
\begin{itemize}
\item Zahed et al. (2021), \emph{The neurophysiology of sleep in Parkinson's disease}: sleep mechanisms (ref.~1, PDF p.~10).
\item Yin et al. (2023), \emph{Pallidal activities during sleep and sleep decoding in dystonia, Huntington's, and Parkinson's disease}: decoding context (ref.~25, PDF p.~11).
\end{itemize}

\end{document}
